% Options for packages loaded elsewhere
\PassOptionsToPackage{unicode}{hyperref}
\PassOptionsToPackage{hyphens}{url}
\documentclass[
]{article}
\usepackage{xcolor}
\usepackage[margin=1in]{geometry}
\usepackage{amsmath,amssymb}
\setcounter{secnumdepth}{-\maxdimen} % remove section numbering
\usepackage{iftex}
\ifPDFTeX
  \usepackage[T1]{fontenc}
  \usepackage[utf8]{inputenc}
  \usepackage{textcomp} % provide euro and other symbols
\else % if luatex or xetex
  \usepackage{unicode-math} % this also loads fontspec
  \defaultfontfeatures{Scale=MatchLowercase}
  \defaultfontfeatures[\rmfamily]{Ligatures=TeX,Scale=1}
\fi
\usepackage{lmodern}
\ifPDFTeX\else
  % xetex/luatex font selection
\fi
% Use upquote if available, for straight quotes in verbatim environments
\IfFileExists{upquote.sty}{\usepackage{upquote}}{}
\IfFileExists{microtype.sty}{% use microtype if available
  \usepackage[]{microtype}
  \UseMicrotypeSet[protrusion]{basicmath} % disable protrusion for tt fonts
}{}
\makeatletter
\@ifundefined{KOMAClassName}{% if non-KOMA class
  \IfFileExists{parskip.sty}{%
    \usepackage{parskip}
  }{% else
    \setlength{\parindent}{0pt}
    \setlength{\parskip}{6pt plus 2pt minus 1pt}}
}{% if KOMA class
  \KOMAoptions{parskip=half}}
\makeatother
\usepackage{color}
\usepackage{fancyvrb}
\newcommand{\VerbBar}{|}
\newcommand{\VERB}{\Verb[commandchars=\\\{\}]}
\DefineVerbatimEnvironment{Highlighting}{Verbatim}{commandchars=\\\{\}}
% Add ',fontsize=\small' for more characters per line
\usepackage{framed}
\definecolor{shadecolor}{RGB}{248,248,248}
\newenvironment{Shaded}{\begin{snugshade}}{\end{snugshade}}
\newcommand{\AlertTok}[1]{\textcolor[rgb]{0.94,0.16,0.16}{#1}}
\newcommand{\AnnotationTok}[1]{\textcolor[rgb]{0.56,0.35,0.01}{\textbf{\textit{#1}}}}
\newcommand{\AttributeTok}[1]{\textcolor[rgb]{0.13,0.29,0.53}{#1}}
\newcommand{\BaseNTok}[1]{\textcolor[rgb]{0.00,0.00,0.81}{#1}}
\newcommand{\BuiltInTok}[1]{#1}
\newcommand{\CharTok}[1]{\textcolor[rgb]{0.31,0.60,0.02}{#1}}
\newcommand{\CommentTok}[1]{\textcolor[rgb]{0.56,0.35,0.01}{\textit{#1}}}
\newcommand{\CommentVarTok}[1]{\textcolor[rgb]{0.56,0.35,0.01}{\textbf{\textit{#1}}}}
\newcommand{\ConstantTok}[1]{\textcolor[rgb]{0.56,0.35,0.01}{#1}}
\newcommand{\ControlFlowTok}[1]{\textcolor[rgb]{0.13,0.29,0.53}{\textbf{#1}}}
\newcommand{\DataTypeTok}[1]{\textcolor[rgb]{0.13,0.29,0.53}{#1}}
\newcommand{\DecValTok}[1]{\textcolor[rgb]{0.00,0.00,0.81}{#1}}
\newcommand{\DocumentationTok}[1]{\textcolor[rgb]{0.56,0.35,0.01}{\textbf{\textit{#1}}}}
\newcommand{\ErrorTok}[1]{\textcolor[rgb]{0.64,0.00,0.00}{\textbf{#1}}}
\newcommand{\ExtensionTok}[1]{#1}
\newcommand{\FloatTok}[1]{\textcolor[rgb]{0.00,0.00,0.81}{#1}}
\newcommand{\FunctionTok}[1]{\textcolor[rgb]{0.13,0.29,0.53}{\textbf{#1}}}
\newcommand{\ImportTok}[1]{#1}
\newcommand{\InformationTok}[1]{\textcolor[rgb]{0.56,0.35,0.01}{\textbf{\textit{#1}}}}
\newcommand{\KeywordTok}[1]{\textcolor[rgb]{0.13,0.29,0.53}{\textbf{#1}}}
\newcommand{\NormalTok}[1]{#1}
\newcommand{\OperatorTok}[1]{\textcolor[rgb]{0.81,0.36,0.00}{\textbf{#1}}}
\newcommand{\OtherTok}[1]{\textcolor[rgb]{0.56,0.35,0.01}{#1}}
\newcommand{\PreprocessorTok}[1]{\textcolor[rgb]{0.56,0.35,0.01}{\textit{#1}}}
\newcommand{\RegionMarkerTok}[1]{#1}
\newcommand{\SpecialCharTok}[1]{\textcolor[rgb]{0.81,0.36,0.00}{\textbf{#1}}}
\newcommand{\SpecialStringTok}[1]{\textcolor[rgb]{0.31,0.60,0.02}{#1}}
\newcommand{\StringTok}[1]{\textcolor[rgb]{0.31,0.60,0.02}{#1}}
\newcommand{\VariableTok}[1]{\textcolor[rgb]{0.00,0.00,0.00}{#1}}
\newcommand{\VerbatimStringTok}[1]{\textcolor[rgb]{0.31,0.60,0.02}{#1}}
\newcommand{\WarningTok}[1]{\textcolor[rgb]{0.56,0.35,0.01}{\textbf{\textit{#1}}}}
\usepackage{graphicx}
\makeatletter
\newsavebox\pandoc@box
\newcommand*\pandocbounded[1]{% scales image to fit in text height/width
  \sbox\pandoc@box{#1}%
  \Gscale@div\@tempa{\textheight}{\dimexpr\ht\pandoc@box+\dp\pandoc@box\relax}%
  \Gscale@div\@tempb{\linewidth}{\wd\pandoc@box}%
  \ifdim\@tempb\p@<\@tempa\p@\let\@tempa\@tempb\fi% select the smaller of both
  \ifdim\@tempa\p@<\p@\scalebox{\@tempa}{\usebox\pandoc@box}%
  \else\usebox{\pandoc@box}%
  \fi%
}
% Set default figure placement to htbp
\def\fps@figure{htbp}
\makeatother
\setlength{\emergencystretch}{3em} % prevent overfull lines
\providecommand{\tightlist}{%
  \setlength{\itemsep}{0pt}\setlength{\parskip}{0pt}}
\usepackage{bookmark}
\IfFileExists{xurl.sty}{\usepackage{xurl}}{} % add URL line breaks if available
\urlstyle{same}
\hypersetup{
  pdftitle={Visualization IsoQuant output},
  hidelinks,
  pdfcreator={LaTeX via pandoc}}

\title{Visualization IsoQuant output}
\author{}
\date{\vspace{-2.5em}Compiled: December 16, 2025}

\begin{document}
\maketitle

\begin{Shaded}
\begin{Highlighting}[]
\FunctionTok{library}\NormalTok{(reticulate)}
\FunctionTok{library}\NormalTok{(glue)}

\FunctionTok{source}\NormalTok{(}\StringTok{"local\_config.R"}\NormalTok{)  }\CommentTok{\# where the path for conda\_path and viewer\_code\_path are set alternatively, set them as below}
\CommentTok{\# conda\_path \textless{}{-} \textquotesingle{}../miniconda3/envs/isoquant\_viewer\_R/\textquotesingle{} \# Replace the path to your conda envirnment path}
\CommentTok{\#viewer\_code\_path \textless{}{-} \textquotesingle{}../IsoQuant\_pipeline/utilities/adapter\_isoquantviewer.R\textquotesingle{} \# Replace the path to your location of the R code}

\FunctionTok{use\_condaenv}\NormalTok{(conda\_path, }\AttributeTok{required =} \ConstantTok{TRUE}\NormalTok{) }
\FunctionTok{use\_python}\NormalTok{(}\FunctionTok{glue}\NormalTok{(}\StringTok{"\{conda\_path\}/bin/python"}\NormalTok{), }\AttributeTok{required =} \ConstantTok{TRUE}\NormalTok{)}
\FunctionTok{source}\NormalTok{(viewer\_code\_path) }
\FunctionTok{py\_config}\NormalTok{() }\CommentTok{\# for checking if the session is configured with the environment you specified}
\end{Highlighting}
\end{Shaded}

\begin{verbatim}
## python:         /Users/kawai/miniconda3/envs/isoquant_viewer_R/bin/python
## libpython:      /Users/kawai/miniconda3/envs/isoquant_viewer_R/lib/libpython3.10.dylib
## pythonhome:     /Users/kawai/miniconda3/envs/isoquant_viewer_R:/Users/kawai/miniconda3/envs/isoquant_viewer_R
## version:        3.10.13 (main, Sep 11 2023, 08:16:02) [Clang 14.0.6 ]
## numpy:          /Users/kawai/miniconda3/envs/isoquant_viewer_R/lib/python3.10/site-packages/numpy
## numpy_version:  2.2.5
## sys:            [builtin module]
## 
## NOTE: Python version was forced by use_python() function
\end{verbatim}

\begin{Shaded}
\begin{Highlighting}[]
\CommentTok{\#output\_directory \textless{}{-} "../prefix/results/isoquant\_output/" \# path to IsoQuantOutput directory}
\CommentTok{\#prefix \textless{}{-} "prefix" \# Name of the project}
\NormalTok{reference\_gtf }\OtherTok{\textless{}{-}} \StringTok{\textquotesingle{}\textasciitilde{}/Documents/ref/geneannotations.gff3\textquotesingle{}} \CommentTok{\# path to reference gene annotation }
\NormalTok{viewer}\OtherTok{\textless{}{-}}\FunctionTok{isoquantViewer}\NormalTok{(}\AttributeTok{output\_directory =}\NormalTok{ output\_directory,        }
    \AttributeTok{prefix =}\NormalTok{ prefix, }
    \AttributeTok{reference\_gtf =}\NormalTok{ reference\_gtf,}
    \AttributeTok{plot\_output =} \StringTok{\textquotesingle{}\textasciitilde{}/Desktop/\textquotesingle{}}
\NormalTok{      )}

\CommentTok{\# * Plots are automatically saved in the folder "./plot\_output" if plot\_output is not defined}
\end{Highlighting}
\end{Shaded}

\paragraph{Visualizing read assignment and
counts}\label{visualizing-read-assignment-and-counts}

Function:

\begin{itemize}
\item
  plot\_pie\_assignment

  Plot pie chart of read assignment (.read\_assignments.tsv.gz; stored
  as .reads\_assignment) from IsoQuant. ``feature\_to\_plot'' are:
  ``assignment\_type'', ``gene\_assignment'', ``Classification''
\item
  plot\_count\_bar

  A bar plot of counts for either specified genes or isoforms. If
  providing gene\_list, plots stacked barplot of isoforms per gene. If
  providing isoform\_list, plots a barplot of specified isoforms.

  Can either be from the raw count or TPM matrix using ``layer'', and
  based on the transcript model or the reference annotation
  (use\_transcript\_model=True or False, default is False). Default is
  plotting all available samples, but can be limited to specific sample
  using ``sample\_id'' = {[}sample1, sample2, and etc.{]}
\end{itemize}

\begin{Shaded}
\begin{Highlighting}[]
\NormalTok{fname}\OtherTok{\textless{}{-}}\FunctionTok{plot\_count\_bar}\NormalTok{(viewer,}\AttributeTok{layer=}\StringTok{\textquotesingle{}count\textquotesingle{}}\NormalTok{,}
          \AttributeTok{isoform\_list =} \FunctionTok{list}\NormalTok{(}\StringTok{\textquotesingle{}ENST00000008440\textquotesingle{}}\NormalTok{,}\StringTok{\textquotesingle{}transcript92.Y.nnic\textquotesingle{}}\NormalTok{),}
          \AttributeTok{use\_transcript\_model=}\ConstantTok{TRUE}\NormalTok{)}
\NormalTok{knitr}\SpecialCharTok{::}\FunctionTok{include\_graphics}\NormalTok{(}\FunctionTok{as.character}\NormalTok{(fname))}
\end{Highlighting}
\end{Shaded}


\includegraphics[width=9.69in]{../../../../../../Desktop/Bar_isoforms_ENST00000008440_transcript92.Y.nnic_TranscriptModel_True}

\begin{Shaded}
\begin{Highlighting}[]
\NormalTok{fname}\OtherTok{\textless{}{-}}\FunctionTok{plot\_pie\_assignment}\NormalTok{(viewer,}\AttributeTok{feature\_to\_plot =} \StringTok{\textquotesingle{}assignment\_type\textquotesingle{}}\NormalTok{)}
\NormalTok{knitr}\SpecialCharTok{::}\FunctionTok{include\_graphics}\NormalTok{(}\FunctionTok{as.character}\NormalTok{(fname))}
\end{Highlighting}
\end{Shaded}

\includegraphics[width=11.53in]{../../../../../../Desktop/Pie_assignment_type}

\paragraph{Visualizing transcripts within a genomic
region}\label{visualizing-transcripts-within-a-genomic-region}

\begin{itemize}
\item
  plot\_transcript\_map

  Visualize transcript(s) for specified gene(s)/ensembl ID within the
  corresponding genomic region(s). Can be based on the transcript model
  or the reference genome. (Based on IsoQuant:plot\_transcript\_map)
\end{itemize}

\begin{Shaded}
\begin{Highlighting}[]
\NormalTok{fname}\OtherTok{\textless{}{-}}\FunctionTok{plot\_transcript\_map}\NormalTok{(viewer,}
                    \AttributeTok{gene\_names=}\FunctionTok{list}\NormalTok{(}\StringTok{"TSPAN6"}\NormalTok{),}
                    \AttributeTok{use\_transcript\_model=}\ConstantTok{TRUE}\NormalTok{)}
\NormalTok{knitr}\SpecialCharTok{::}\FunctionTok{include\_graphics}\NormalTok{(}\FunctionTok{as.character}\NormalTok{(fname))}
\end{Highlighting}
\end{Shaded}

\includegraphics[width=15.5in]{../../../../../../Desktop/TranscriptMap_ENSG00000000003_TranscriptModel_True}

\begin{itemize}
\item
  plot\_genomic\_region

  Visualize a genomic region defined by setting either:

  \begin{itemize}
  \tightlist
  \item
    region = `X:100627109-100639991'
  \item
    gene\_name = `TSPAN6'
  \item
    or Ensembl\_ID = `ENSG00000000003'
  \end{itemize}

  using pyGenomeTracks classes
  (\url{https://pygenometracks.readthedocs.io/en/latest/}).

  The default ordered tracks being plotted are
  {[}`ref',`model',`reads'{]},can be set using ``plot\_tracks''.

  \begin{itemize}
  \tightlist
  \item
    `ref': Reference GTF
  \item
    `model': IsoQuant transcript model GTF
  \item
    `reads': IsoQuant corrected alignment BED
  \end{itemize}

  The default parameters for each track can be overwritten using
  tracks\_params, by setting key-value pairs for any supported track
  property. (See the list of all available options at at
  \url{https://pygenometracks.readthedocs.io/en/latest/content/all_tracks.html})
\end{itemize}

\begin{Shaded}
\begin{Highlighting}[]
\NormalTok{fname}\OtherTok{\textless{}{-}}\FunctionTok{plot\_genomic\_region}\NormalTok{(viewer,}
                    \AttributeTok{gene\_name =} \StringTok{\textquotesingle{}TSPAN6\textquotesingle{}}\NormalTok{)}
\NormalTok{knitr}\SpecialCharTok{::}\FunctionTok{include\_graphics}\NormalTok{(}\FunctionTok{as.character}\NormalTok{(fname))}
\end{Highlighting}
\end{Shaded}

\includegraphics[width=16.44in,height=1\textheight]{../../../../../../Desktop/GenomeTrack_X_100627109-100639991}

\end{document}
